\section{Bảng lệnh Git tra cứu nhanh}
\label{sec:pa-banglenh}

Phụ lục này tổng hợp các lệnh Git dùng thường xuyên nhất trong học phần ASPX,
sắp xếp theo nhóm công việc từ khởi tạo kho chứa đến thao tác nâng cao.
Mỗi bảng gồm bốn cột: nhóm công việc, lệnh, mục đích và ví dụ chạy được.
Sinh viên nên đọc kèm Chương 2 và Chương 3 của báo cáo để hiểu ngữ cảnh sử dụng,
cùng tài liệu chính thức của Git \cite{progit} và hướng dẫn tham khảo của Atlassian \cite{atlassiangit}.

Hai ví dụ tổng hợp dưới đây minh họa luồng làm việc cơ bản và luồng làm việc theo nhánh
với project Web Forms. Mọi lệnh đều chạy được trên Git Bash hoặc PowerShell đã cài Git.

\begin{lstlisting}[language=bash]
# Basic workflow: config once, then add and commit
git config --global user.name "Nguyen Van A"
git config --global user.email "nguyenvana@example.com"
git init AspxDemo_NET481
cd AspxDemo_NET481
git status
git add .
git commit -m "feat: khoi tao project ASPX"
git log --oneline -n 5
\end{lstlisting}

\begin{lstlisting}[language=bash]
# Branch workflow: feature branch, merge, push to remote
git checkout -b feature/login-form
git add Default.aspx Default.aspx.cs
git commit -m "feat: them form dang nhap"
git checkout main
git merge feature/login-form
git remote add origin https://github.com/MinhTriTM/AspxDemo_NET481.git
git push -u origin main
\end{lstlisting}

\subsection{Nhóm Khởi tạo}
Tạo kho chứa mới hoặc lấy kho chứa có sẵn về máy trước khi bắt đầu mọi công việc khác.
\begin{table}[H]
\centering
\small
\begin{tabular}{p{1.7cm} p{3.2cm} p{4.5cm} p{4.6cm}}
\toprule
Nhóm & Lệnh & Mục đích & Ví dụ \\
\midrule
Khởi tạo & \verb|git init| & Tạo kho Git mới trong thư mục hiện tại & \verb|git init AspxDemo_NET481| \\
Khởi tạo & \verb|git clone <url>| & Sao chép toàn bộ kho chứa từ xa về máy & \verb|git clone --depth 1 <url>| \\
Khởi tạo & \verb|git config user.name| & Khai báo tên tác giả cho commit & \verb|git config user.name "An"| \\
Khởi tạo & \verb|git config user.email| & Khai báo email tác giả cho commit & \verb|git config user.email a@b.vn| \\
Khởi tạo & \verb|git config --list| & Liệt kê toàn bộ cấu hình đang áp dụng & \verb|git config --global --list| \\
\bottomrule
\end{tabular}
\caption{Bảng lệnh nhóm Khởi tạo}
\label{tab:pa-khoitao}
\end{table}

\subsection{Nhóm Commit}
Ghi lại thay đổi theo từng bước nhỏ với thông điệp rõ ràng theo Conventional Commits.
\begin{table}[H]
\centering
\small
\begin{tabular}{p{1.7cm} p{3.2cm} p{4.5cm} p{4.6cm}}
\toprule
Nhóm & Lệnh & Mục đích & Ví dụ \\
\midrule
Commit & \verb|git status| & Xem trạng thái Working và Staging & \verb|git status -s| \\
Commit & \verb|git add <file>| & Đưa một file vào vùng Staging & \verb|git add Default.aspx| \\
Commit & \verb|git add .| & Đưa mọi thay đổi vào vùng Staging & \verb|git add -A| \\
Commit & \verb|git commit -m| & Chụp ảnh Staging thành một commit & \verb|git commit -m "feat: login"| \\
Commit & \verb|git commit --amend| & Sửa commit gần nhất thay vì tạo mới & \verb|git commit --amend --no-edit| \\
\bottomrule
\end{tabular}
\caption{Bảng lệnh nhóm Commit}
\label{tab:pa-commit}
\end{table}

\subsection{Nhóm Lịch sử}
Xem lại quá khứ của kho chứa để hiểu ai đã sửa gì và so sánh các phiên bản.
\begin{table}[H]
\centering
\small
\begin{tabular}{p{1.7cm} p{3.2cm} p{4.5cm} p{4.6cm}}
\toprule
Nhóm & Lệnh & Mục đích & Ví dụ \\
\midrule
Lịch sử & \verb|git log| & Xem lịch sử commit dạng đầy đủ & \verb|git log -n 5| \\
Lịch sử & \verb|git log --oneline| & Xem lịch sử gọn một dòng một commit & \verb|git log --oneline --graph| \\
Lịch sử & \verb|git diff| & So sánh Working với Staging & \verb|git diff Default.aspx| \\
Lịch sử & \verb|git diff --staged| & So sánh Staging với commit gần nhất & \verb|git diff --cached| \\
Lịch sử & \verb|git show <commit>| & Xem chi tiết một commit cụ thể & \verb|git show HEAD~1| \\
\bottomrule
\end{tabular}
\caption{Bảng lệnh nhóm Lịch sử}
\label{tab:pa-lichsu}
\end{table}

\subsection{Nhóm Nhánh}
Mỗi tính năng làm trên một nhánh riêng để tránh đè code khi ba người cùng sửa project.
\begin{table}[H]
\centering
\small
\begin{tabular}{p{1.7cm} p{3.2cm} p{4.5cm} p{4.6cm}}
\toprule
Nhóm & Lệnh & Mục đích & Ví dụ \\
\midrule
Nhánh & \verb|git branch| & Liệt kê các nhánh trong kho chứa & \verb|git branch -a| \\
Nhánh & \verb|git branch <ten>| & Tạo nhánh mới từ HEAD hiện tại & \verb|git branch feature/x| \\
Nhánh & \verb|git switch <nhanh>| & Chuyển sang nhánh khác an toàn & \verb|git switch main| \\
Nhánh & \verb|git switch -c <ten>| & Tạo và chuyển ngay sang nhánh mới & \verb|git switch -c feature/y| \\
Nhánh & \verb|git branch -d <ten>| & Xóa nhánh đã merge xong & \verb|git branch -d feature/x| \\
\bottomrule
\end{tabular}
\caption{Bảng lệnh nhóm Nhánh}
\label{tab:pa-nhanh}
\end{table}

\subsection{Nhóm Merge/Rebase}
Gộp công việc từ các nhánh lại với nhau và xử lý conflict khi hai người sửa cùng một chỗ.
\begin{table}[H]
\centering
\small
\begin{tabular}{p{1.7cm} p{3.2cm} p{4.5cm} p{4.6cm}}
\toprule
Nhóm & Lệnh & Mục đích & Ví dụ \\
\midrule
Merge & \verb|git merge <nhanh>| & Gộp nhánh chỉ định vào nhánh hiện tại & \verb|git merge feature/x| \\
Merge & \verb|git merge --abort| & Hủy bỏ lần merge đang dở dang & \verb|git merge --abort| \\
Rebase & \verb|git rebase <nhanh>| & Đặt lại commit lên đầu nhánh khác & \verb|git rebase main| \\
Rebase & \verb|git rebase -i HEAD~3| & Rebase tương tác để gộp và sửa commit & \verb|git rebase -i HEAD~5| \\
Rebase & \verb|git cherry-pick <c>| & Sao chép một commit sang nhánh hiện tại & \verb|git cherry-pick a1b2c3d| \\
\bottomrule
\end{tabular}
\caption{Bảng lệnh nhóm Merge/Rebase}
\label{tab:pa-merge}
\end{table}

\subsection{Nhóm Hoàn tác}
Sửa sai an toàn: ưu tiên các lệnh tạo commit mới thay vì xóa lịch sử đã chia sẻ.
\begin{table}[H]
\centering
\small
\begin{tabular}{p{1.7cm} p{3.2cm} p{4.5cm} p{4.6cm}}
\toprule
Nhóm & Lệnh & Mục đích & Ví dụ \\
\midrule
Hoàn tác & \verb|git restore <file>| & Hủy sửa trong Working về bản commit & \verb|git restore Web.config| \\
Hoàn tác & \verb|git restore --staged| & Gỡ file khỏi Staging giữ nguyên sửa & \verb|git restore --staged .| \\
Hoàn tác & \verb|git reset --soft HEAD~1| & Lùi một commit giữ Staging nguyên & \verb|git reset --soft HEAD~2| \\
Hoàn tác & \verb|git reset --hard HEAD| & Xóa mọi sửa chưa commit về HEAD & \verb|git reset --hard origin/main| \\
Hoàn tác & \verb|git revert <commit>| & Tạo commit mới đảo ngược commit cũ & \verb|git revert HEAD| \\
\bottomrule
\end{tabular}
\caption{Bảng lệnh nhóm Hoàn tác}
\label{tab:pa-hoantac}
\end{table}

\subsection{Nhóm Stash/Tag}
Cất tạm công việc đang dở và đánh dấu các mốc phiên bản quan trọng của đồ án.
\begin{table}[H]
\centering
\small
\begin{tabular}{p{1.7cm} p{3.2cm} p{4.5cm} p{4.6cm}}
\toprule
Nhóm & Lệnh & Mục đích & Ví dụ \\
\midrule
Stash & \verb|git stash push -m| & Cất tạm sửa chưa commit kèm ghi chú & \verb|git stash push -m "wip"| \\
Stash & \verb|git stash list| & Liệt kê các bản cất tạm đang giữ & \verb|git stash list| \\
Stash & \verb|git stash pop| & Khôi phục bản cất tạm gần nhất & \verb|git stash pop stash@{0}| \\
Tag & \verb|git tag <ten>| & Gắn thẻ phiên bản tại commit hiện tại & \verb|git tag v1.0| \\
Tag & \verb|git tag -a <ten>| & Gắn thẻ có chú thích cho bản phát hành & \verb|git tag -a v1.0 -m "release"| \\
\bottomrule
\end{tabular}
\caption{Bảng lệnh nhóm Stash/Tag}
\label{tab:pa-stash}
\end{table}

\subsection{Nhóm Remote}
Đồng bộ kho chứa nội bộ với GitHub và GitLab để cả nhóm cùng thấy code mới nhất.
\begin{table}[H]
\centering
\small
\begin{tabular}{p{1.7cm} p{3.2cm} p{4.5cm} p{4.6cm}}
\toprule
Nhóm & Lệnh & Mục đích & Ví dụ \\
\midrule
Remote & \verb|git remote -v| & Liệt kê các kho chứa từ xa & \verb|git remote -v| \\
Remote & \verb|git remote add origin| & Gắn kho chứa từ xa tên origin & \verb|git remote add upstream <url>| \\
Remote & \verb|git fetch| & Tải commit mới về chưa gộp & \verb|git fetch origin| \\
Remote & \verb|git pull| & Tải về và gộp nhánh từ xa & \verb|git pull --rebase| \\
Remote & \verb|git push| & Đẩy commit lên kho chứa từ xa & \verb|git push -u origin main| \\
\bottomrule
\end{tabular}
\caption{Bảng lệnh nhóm Remote}
\label{tab:pa-remote}
\end{table}

\subsection{Nhóm Nâng cao}
Các lệnh cứu hộ và dọn dẹp kho chứa khi gặp sự cố hoặc cần kiểm tra lịch sử sâu.
\begin{table}[H]
\centering
\small
\begin{tabular}{p{1.7cm} p{3.2cm} p{4.5cm} p{4.6cm}}
\toprule
Nhóm & Lệnh & Mục đích & Ví dụ \\
\midrule
Nâng cao & \verb|git reflog| & Xem lịch sử di chuyển của HEAD & \verb|git reflog -n 10| \\
Nâng cao & \verb|git bisect start| & Bắt đầu tìm commit gây lỗi & \verb|git bisect bad HEAD| \\
Nâng cao & \verb|git clean -fd| & Xóa file và thư mục chưa theo dõi & \verb|git clean -fdn| \\
Nâng cao & \verb|git rm --cached| & Ngừng theo dõi file giữ lại trên đĩa & \verb|git rm --cached app.log| \\
Nâng cao & \verb|git worktree list| & Liệt kê các cây làm việc đang dùng & \verb|git worktree add ../hotfix| \\
\bottomrule
\end{tabular}
\caption{Bảng lệnh nhóm Nâng cao}
\label{tab:pa-nangcao}
\end{table}

Tổng cộng phụ lục này liệt kê 45 lệnh thuộc 9 nhóm công việc, đủ cho sinh viên
tra cứu nhanh trong suốt quá trình làm đồ án ASPX theo nhóm ba người.
